\documentclass[11pt,a4paper]{article}
\usepackage[utf8]{inputenc}
\usepackage[T1]{fontenc}

\usepackage[spanish]{babel}
\usepackage[scale=0.86,top=1.4cm,bottom=1.4cm]{geometry}
\usepackage{enumitem}
\usepackage{titlesec}
\usepackage{xcolor}
\usepackage[hidelinks]{hyperref}
\usepackage{array}

\definecolor{accent}{RGB}{20,20,20}
\definecolor{rule}{RGB}{120,120,120}

\pagestyle{empty}
\setlength{\parindent}{0pt}

% Section styling
\titleformat{\section}
  {\large\bfseries\scshape}
  {}{0pt}{}[\vspace{-0.6ex}\color{rule}\titlerule]
\titlespacing{\section}{0pt}{2.2ex}{1.2ex}

\setlist[itemize]{leftmargin=1.2em,itemsep=0.15ex,topsep=0.4ex,parsep=0pt}

% Entry command: {role}{org}{dates}{location}
\newcommand{\entry}[4]{%
  \noindent\begin{tabular*}{\textwidth}{@{}l@{\extracolsep{\fill}}r@{}}
  \textbf{#1} & \textbf{#4} \\
  \textit{#2} & \textit{#3} \\
  \end{tabular*}\vspace{0.3ex}%
}

\begin{document}

% ===== Header =====
\begin{center}
  {\LARGE\bfseries Darwin Cristhian Turpo Quispe}\\[0.4ex]
  {\large Desarrollador Full Stack Junior \ $\cdot$\ Python $\cdot$ FastAPI $\cdot$ Django $\cdot$ JavaScript}\\[0.9ex]
  +51 957565165 \quad$\vert$\quad darwin.turpo17@gmail.com \quad$\vert$\quad Arequipa, PE\\[0.4ex]
  \href{https://linkedin.com/in/darwin-turpo-813083313}{LinkedIn}
  \quad$\vert$\quad
  \href{https://github.com/darwintq}{GitHub}
  \quad$\vert$\quad
  \href{https://portfolio-ai-cyan-seven.vercel.app/}{Portafolio}
\end{center}

% ===== Perfil =====
\section{Perfil Profesional}
Estudiante de Ingeniería de Software con IA e Ingeniería de Telecomunicaciones, orientado al desarrollo backend y la integración de sistemas. Construyo APIs y arquitecturas modulares con Python, FastAPI y Django, y trabajo cómodo con SQL, Git y entornos Linux. Enfocado en entregar software bien estructurado, mantenible y probado, aportando valor técnico desde el primer día.

% ===== Experiencia =====
\section{Experiencia y Proyectos}
\entry{Desarrollador Full Stack}{Pontificia Universidad Católica del Perú (PUCP)}{Ene 2026 -- Ago 2026}{Lima, PE}
\begin{itemize}
  \item Diseñé una arquitectura modular por capas con FastAPI y WebSockets que desacopló 3 subsistemas del CubeSat 1U, permitiendo pruebas e integración independientes por subsistema.
  \item Centralicé el registro de más de 5 sensores en una base de datos relacional y construí dashboards de monitoreo en Django, habilitando la trazabilidad de telemetría en tiempo real.
  \item Implementé una estrategia de branching y revisión de código en Git para un equipo de 4 personas, reduciendo de forma notable los conflictos de integración.
\end{itemize}

% ===== Formación Académica =====
\section{Formación Académica}
\entry{Ingeniería de Software con Inteligencia Artificial}{SENATI}{2024 -- actualidad}{Perú}
\vspace{0.6ex}
\entry{Ingeniería de Telecomunicaciones}{Universidad Nacional de San Agustín de Arequipa}{2023 -- actualidad}{Arequipa, PE}

% ===== Competencias =====
\section{Competencias Técnicas}
\begin{itemize}
  \item \textbf{Lenguajes:} Python, JavaScript, TypeScript, Java, PHP
  \item \textbf{Backend / APIs:} FastAPI, Django, Laravel, REST APIs, WebSockets
  \item \textbf{Bases de Datos:} Oracle SQL, MySQL
  \item \textbf{Frontend:} HTML5, CSS3
  \item \textbf{Herramientas y Sistemas:} Git, Linux (Red Hat), Android Studio, Flutter, DrawIO
  \item \textbf{Productividad con IA:} Claude Code, OpenCode, ChatGPT, Gemini
\end{itemize}

% ===== Formación Complementaria =====
\section{Formación Complementaria}
\begin{itemize}
  \item Python Essentials 2 --- Cisco Academy --- 160 horas --- 2024
  \item SQL Database --- Oracle Academy --- 40 horas --- 2025
  \item IT Essentials --- INFOUNSA --- 20 horas --- 2025
  \item CADE Universitario 2026 --- IPAE --- 3 días (liderazgo y emprendimiento)
\end{itemize}

% ===== Habilidades Blandas =====
\section{Habilidades Blandas}
Trabajo en equipo, aprendizaje continuo, pensamiento analítico, adaptabilidad, compromiso y puntualidad.

% ===== Idiomas =====
\section{Idiomas}
\begin{itemize}
  \item \textbf{Español:} Nativo
  \item \textbf{Inglés:} Intermedio --- Cultural Arequipa (Enero 2025 -- Actualidad)
\end{itemize}

\end{document}